\section{Introduction}
\label{sec:introduction}

Teaching the basic behavior of quantum circuits is difficult in part because the underlying
phenomena -- a qubit existing in superposition rather than a definite 0 or 1, a measurement
irreversibly collapsing that superposition to a single classical outcome, two qubits becoming
entangled such that neither can be described independently of the other -- have no
counterpart at all in classical, everyday intuition, unlike an RC charging curve or a
mechanical oscillator, each of which at least resembles something a student has already
observed. The conventional response pairs the mathematics of complex amplitudes and unitary
matrices with a notebook-based simulator or an online circuit composer such as IBM's Quantum
Composer or Google's Cirq, in which a student drags gates onto a schematic timeline and reads
off a resulting statevector or histogram. These tools are effective and widely used, but each
presents the circuit as a two-dimensional diagram manipulated with a mouse, once removed from
the kind of first-person, spatial construction a sandbox game affords.

Sandbox games occupy a distinctive position relative to this approach. Minecraft in
particular has an established record as an informal STEM learning environment
(Section~\ref{sec:related-work}), in part because constructing something inside it requires
systems thinking comparable to an engineering task while remaining playful and low
stakes~\cite{nebel2016mining}. Minecraft's own redstone subsystem is already a widely used
entry point into classical digital logic, examined further in
Section~\ref{sec:minecraft-mechanics}. No equivalent exists, however, for \emph{quantum}
circuits within the base game: redstone signals are boolean and locally deterministic, with
no built-in notion of superposition, complex amplitude, or entanglement between distant
signals.

This paper describes QuantumCraft~\cite{picos2026quantumcraft}, a modification (``mod'') for
the Fabric mod loader, released as free and open-source software under the MIT license, that
adds a small, complete set of quantum-circuit components to Minecraft: a qubit-line block
establishing one of up to three parallel timeline rows, a single cyclable gate block offering
the Pauli~$X$, $Y$, $Z$, Hadamard, $S$, and $T$ gates, a pair of control and target marker
blocks that together represent a controlled-NOT or, with two controls, a Toffoli gate, and a
handheld probe reading out either the exact probability of every basis state or a genuine
collapsing measurement. Placing gate and marker blocks along a qubit line, left to right in
time order, constructs an actual quantum circuit within the game world; probing that circuit
runs a real complex-amplitude statevector simulator, applying the exact unitary matrix for
each placed gate in sequence, rather than a scripted approximation of what a quantum computer
would report. An entangled Bell state built this way genuinely has zero probability of
appearing in the $|01\rangle$ or $|10\rangle$ basis states, and repeated measurements of it
genuinely reproduce a 50/50 split between $|00\rangle$ and $|11\rangle$, in the same sense
that CircuitSimCraft's~\cite{picos2026circuitsimcraft} RC charging curve is a real solution to the
underlying differential equation rather than an animation tuned to look right.

The remainder of this paper is organized as follows. Section~\ref{sec:related-work}
situates QuantumCraft relative to existing quantum-circuit simulators, the foundational
quantum-computing literature, and the literature on Minecraft in education.
Section~\ref{sec:minecraft-mechanics} examines which of Minecraft's own construction
mechanics the mod builds upon, and why redstone's existing role as an entry point into
classical digital logic does not, by itself, extend to representing a qubit.
Section~\ref{sec:architecture} describes the mod's own software architecture, in particular
the algorithm used to translate a player's in-world gate placement into a quantum circuit.
Section~\ref{sec:quantum-simulator} describes the statevector simulator itself, including the
gate set, the general controlled-gate and Toffoli mechanisms, and Born-rule measurement.
Section~\ref{sec:pedagogy} discusses the pedagogical rationale underlying each design choice,
and Section~\ref{sec:experiments} gives five complete, reproducible worked circuits -- a
single-qubit superposition, a Bell state, a phase-kickback demonstration, a three-qubit GHZ
state, and a Toffoli-gate truth-table verification -- built directly into the same world,
with exact gate sequences and closed-form expected results for each.
Section~\ref{sec:validation} describes how the simulator was verified during development.
Section~\ref{sec:limitations} discusses current limitations and planned classroom
evaluation, and Section~\ref{sec:conclusion} concludes.
